\documentclass[10pt,a4paper]{amsart}
\usepackage{amsmath,amssymb,amsthm,graphicx}
\usepackage{verbatim}
\newenvironment{reference}{\comment}{\endcomment}
%\newenvironment{reference}{}{}
\newenvironment{slogan}{\comment}{\endcomment}
\newenvironment{history}{\comment}{\endcomment}

% For commutative diagrams we use Xy-pic
\usepackage[all]{xy}

% We use 2cell for 2-commutative diagrams.
\xyoption{2cell}
% Theorem environments.
%
\theoremstyle{plain}
\newtheorem{theorem}[subsection]{Theorem}
\newtheorem{proposition}[subsection]{Proposition}
\newtheorem{lemma}[subsection]{Lemma}

\theoremstyle{definition}
\newtheorem{definition}[subsection]{Definition}
\newtheorem{example}[subsection]{Example}
\newtheorem{exercise}[subsection]{Exercise}
\newtheorem{situation}[subsection]{Situation}

\theoremstyle{remark}
\newtheorem{remark}[subsection]{Remark}
\newtheorem{remarks}[subsection]{Remarks}

\numberwithin{equation}{subsection}

% Macros
%
\def\lim{\mathop{\mathrm{lim}}\nolimits}
\def\colim{\mathop{\mathrm{colim}}\nolimits}
\def\Spec{\mathop{\mathrm{Spec}}}
\def\Hom{\mathop{\mathrm{Hom}}\nolimits}
\def\Ext{\mathop{\mathrm{Ext}}\nolimits}
\def\SheafHom{\mathop{\mathcal{H}\!\mathit{om}}\nolimits}
\def\SheafExt{\mathop{\mathcal{E}\!\mathit{xt}}\nolimits}
\def\Sch{\mathit{Sch}}
\def\Mor{\mathop{\mathrm{Mor}}\nolimits}
\def\Ob{\mathop{\mathrm{Ob}}\nolimits}
\def\Sh{\mathop{\mathit{Sh}}\nolimits}
\def\NL{\mathop{N\!L}\nolimits}
\def\CH{\mathop{\mathrm{CH}}\nolimits}
\def\proetale{{pro\text{-}\acute{e}tale}}
\def\etale{{\acute{e}tale}}
\def\QCoh{\mathit{QCoh}}
\def\Ker{\mathop{\mathrm{Ker}}}
\def\Im{\mathop{\mathrm{Im}}}
\def\Coker{\mathop{\mathrm{Coker}}}
\def\Coim{\mathop{\mathrm{Coim}}}

% Boxtimes
%
\DeclareMathSymbol{\boxtimes}{\mathbin}{AMSa}{"02}

%
% Macros for moduli stacks/spaces
%
\def\QCohstack{\mathcal{QC}\!\mathit{oh}}
\def\Cohstack{\mathcal{C}\!\mathit{oh}}
\def\Spacesstack{\mathcal{S}\!\mathit{paces}}
\def\Quotfunctor{\mathrm{Quot}}
\def\Hilbfunctor{\mathrm{Hilb}}
\def\Curvesstack{\mathcal{C}\!\mathit{urves}}
\def\Polarizedstack{\mathcal{P}\!\mathit{olarized}}
\def\Complexesstack{\mathcal{C}\!\mathit{omplexes}}
% \Pic is the operator that assigns to X its picard group, usage \Pic(X)
% \Picardstack_{X/B} denotes the Picard stack of X over B
% \Picardfunctor_{X/B} denotes the Picard functor of X over B
\def\Pic{\mathop{\mathrm{Pic}}\nolimits}
\def\Picardstack{\mathcal{P}\!\mathit{ic}}
\def\Picardfunctor{\mathrm{Pic}}
\def\Deformationcategory{\mathcal{D}\!\mathit{ef}}
\usepackage{cleveref}
\crefname{originallemma}{original lemma}{original lemmas}
\Crefname{originallemma}{Original lemma}{Original lemmas}
\makeatletter
\newlabel{lemma-localize-topos}{{04GZ}{source}{}{sourcepointer-1}{}}
\newlabel{lemma-localize-topos@cref}{{[originallemma][0][]04GZ}{[1][0][]source}{}{}{}}
\makeatother
\crefname{originallemma}{original lemma}{original lemmas}
\Crefname{originallemma}{Original lemma}{Original lemmas}
\makeatletter
\newlabel{lemma-topos-good-site}{{03CI}{source}{}{sourcepointer-2}{}}
\newlabel{lemma-topos-good-site@cref}{{[originallemma][0][]03CI}{[1][0][]source}{}{}{}}
\makeatother
\crefname{originallemma}{original lemma}{original lemmas}
\Crefname{originallemma}{Original lemma}{Original lemmas}
\makeatletter
\newlabel{lemma-localize-given-products}{{03CE}{source}{}{sourcepointer-3}{}}
\newlabel{lemma-localize-given-products@cref}{{[originallemma][0][]03CE}{[1][0][]source}{}{}{}}
\makeatother
\crefname{originallemma}{original lemma}{original lemmas}
\Crefname{originallemma}{Original lemma}{Original lemmas}
\makeatletter
\newlabel{categories-lemma-yoneda}{{001P}{source}{}{sourcepointer-4}{}}
\newlabel{categories-lemma-yoneda@cref}{{[originallemma][0][]001P}{[1][0][]source}{}{}{}}
\makeatother
\crefname{originallemma}{original lemma}{original lemmas}
\Crefname{originallemma}{Original lemma}{Original lemmas}
\makeatletter
\newlabel{categories-lemma-characterize-mono-epi}{{08LR}{source}{}{sourcepointer-5}{}}
\newlabel{categories-lemma-characterize-mono-epi@cref}{{[originallemma][0][]08LR}{[1][0][]source}{}{}{}}
\makeatother
\crefname{originallemma}{original lemma}{original lemmas}
\Crefname{originallemma}{Original lemma}{Original lemmas}
\makeatletter
\newlabel{lemma-restrict-back}{{00Y2}{source}{}{sourcepointer-6}{}}
\newlabel{lemma-restrict-back@cref}{{[originallemma][0][]00Y2}{[1][0][]source}{}{}{}}
\makeatother
\title{Subobjects of the final object and subtopoi}
\author{The Stacks Project Authors}
\date{stacks-project:sites.tex:local-snapshot-20261012; GNU FDL 1.2 or later}
\begin{document}
\maketitle
\noindent\textit{Editorial note.} All delivered mathematical text is the author's own, copied from the pinned Stacks Project chapter file sites-cohomology.tex; no statement or proof step was rewritten, reordered or removed.\par\medskip
\noindent\textit{Editorial note.} The proof is the author's complete argument; every step the author writes out is inside the counted span and no closing sentence is a placeholder.\par\medskip
\noindent\textit{Editorial note.} Every labelled reference occurring inside the counted spans is carried in assembly.external\_references as a non-load-bearing navigation pointer (same-chapter ones at their definition site, cross-chapter ones at the invocation inside the counted span), so no load-bearing mathematics is replaced by a stub. No bibliography citation occurs inside the counted spans.\par\medskip
\noindent\textit{Editorial note.} The checked chapter list for same-name collisions is the set of stacks chapters staged in this repository: chow, cohomology, divisors, etale-cohomology, flat, more-algebra, more-morphisms, morphisms, spaces-more-morphisms, varieties together with this card's chapter sites-cohomology.\par\medskip
\noindent\textit{Editorial note.} The preamble is the chapter's own retained minimal preamble (sources/stacks-sites-cohomology/retained-preamble.tex), a verbatim excerpt of the native preamble: the Xy-pic setup and the native theorem declarations with the mathematical macros.\par\medskip
\noindent\textit{Allowed standard prerequisites (editorial).}\begin{itemize}
\item The slice topos and its localization morphism. Same chapter, Lemma lemma-localize-topos. Supplies the identification of the slice of the topos with the localized category of sheaves.
\item Replacing the topos by a good site. Same chapter, Lemma lemma-topos-good-site. Allows the author to assume that the sheaf is representable by an object of a site with subcanonical topology and finite limits.
\item The pushforward along a localization when the site has products. Same chapter, Lemma lemma-localize-given-products and Lemma lemma-restrict-back. Supply the computations of the pushforward and of the composite identity used in the two directions.
\end{itemize}\medskip
\ifdefined\hypertarget\hypertarget{sourcepointer-1}{}\fi
\noindent\textit{Source reference.} Original lemma 04GZ: sources/stacks-sites/pinned-original-source.tex, source lines 10343--10343. Navigation pointer invoked at proof line 10343; the counted span contains the pointer only and the cited result is not reproduced inside it, so the pointer is registered as a non-load-bearing navigation stub in the established form for Stacks material. No bibliography citation is involved.\par\medskip
\ifdefined\hypertarget\hypertarget{sourcepointer-2}{}\fi
\noindent\textit{Source reference.} Original lemma 03CI: sources/stacks-sites/pinned-original-source.tex, source lines 10345--10345. Navigation pointer invoked at proof line 10345; the counted span contains the pointer only and the cited result is not reproduced inside it, so the pointer is registered as a non-load-bearing navigation stub in the established form for Stacks material. No bibliography citation is involved.\par\medskip
\ifdefined\hypertarget\hypertarget{sourcepointer-3}{}\fi
\noindent\textit{Source reference.} Original lemma 03CE: sources/stacks-sites/pinned-original-source.tex, source lines 10362--10362. Navigation pointer invoked at proof line 10362; the counted span contains the pointer only and the cited result is not reproduced inside it, so the pointer is registered as a non-load-bearing navigation stub in the established form for Stacks material. No bibliography citation is involved.\par\medskip
\ifdefined\hypertarget\hypertarget{sourcepointer-4}{}\fi
\noindent\textit{Source reference.} Original lemma 001P: sources/stacks-sites/pinned-original-source.tex, source lines 10368--10368. Navigation pointer invoked at proof line 10368; the counted span contains the pointer only and the cited result is not reproduced inside it, so the pointer is registered as a non-load-bearing navigation stub in the established form for Stacks material. No bibliography citation is involved.\par\medskip
\ifdefined\hypertarget\hypertarget{sourcepointer-5}{}\fi
\noindent\textit{Source reference.} Original lemma 08LR: sources/stacks-sites/pinned-original-source.tex, source lines 10370--10370. Navigation pointer invoked at proof line 10370; the counted span contains the pointer only and the cited result is not reproduced inside it, so the pointer is registered as a non-load-bearing navigation stub in the established form for Stacks material. No bibliography citation is involved.\par\medskip
\ifdefined\hypertarget\hypertarget{sourcepointer-6}{}\fi
\noindent\textit{Source reference.} Original lemma 00Y2: sources/stacks-sites/pinned-original-source.tex, source lines 10375--10375. Navigation pointer invoked at proof line 10375; the counted span contains the pointer only and the cited result is not reproduced inside it, so the pointer is registered as a non-load-bearing navigation stub in the established form for Stacks material. No bibliography citation is involved.\par\medskip
\begin{lemma}
\label{lemma-open-subtopos}
Let $\mathcal{C}$ be a site. Let $\mathcal{F}$ be a sheaf on
$\mathcal{C}$. The following are equivalent
\begin{enumerate}
\item $\mathcal{F}$ is a subobject of the final object of
$\Sh(\mathcal{C})$, and
\item the topos $\Sh(\mathcal{C})/\mathcal{F}$ is a subtopos of
$\Sh(\mathcal{C})$.
\end{enumerate}
\end{lemma}

\begin{proof}
We have seen in Lemma \ref{lemma-localize-topos} that
$\Sh(\mathcal{C})/\mathcal{F}$ is a topos. In fact, we recall the
proof. First we apply Lemma \ref{lemma-topos-good-site}
to see that we may assume $\mathcal{C}$ is a site with a subcanonical
topology, fibre products, a final object $X$, and an object $U$ with
$\mathcal{F} = h_U$. The proof of
Lemma \ref{lemma-localize-topos}
shows that the morphism of topoi
$j_\mathcal{F} : \Sh(\mathcal{C})/\mathcal{F} \to \Sh(\mathcal{C})$
is equal (modulo certain identifications) to the localization morphism
$j_U : \Sh(\mathcal{C}/U) \to \Sh(\mathcal{C})$.

\medskip\noindent
Assume (2). This means that $j_U^{-1}j_{U, *}\mathcal{G} \to \mathcal{G}$
is an isomorphism for all sheaves $\mathcal{G}$ on $\mathcal{C}/U$.
For any object $Z/U$ of $\mathcal{C}/U$ we have
$$
(j_{U, *}h_{Z/U})(U) = \Mor_{\mathcal{C}/U}(U \times_X U/U, Z/U)
$$
by Lemma \ref{lemma-localize-given-products}.
Setting $\mathcal{G} = h_{Z/U}$ in the equality above we obtain
$$
\Mor_{\mathcal{C}/U}(U \times_X U/U, Z/U) = \Mor_{\mathcal{C}/U}(U, Z/U)
$$
for all $Z/U$. By Yoneda's lemma
(Categories, Lemma \ref{categories-lemma-yoneda})
this implies $U \times_X U = U$. By
Categories, Lemma \ref{categories-lemma-characterize-mono-epi}
$U \to X$ is a monomorphism, in other words (1) holds.

\medskip\noindent
Assume (1). Then $j_U^{-1} j_{U, *} = \text{id}$ by
Lemma \ref{lemma-restrict-back}.
\end{proof}

\end{document}
